% G1 theory-only correction and verification contract (2026-11-04 UTC).
% Restricts claims to the linear MQS baseline. Not an assertion that G1 has passed.
\section{G1 基础方程、近似空间与可信参考解的重新核查}
\label{sec:g1-foundation-recheck}

\subsection{首先区分三个不同的数学对象}
G1 最终要近似的是无限均匀导电海水中的连续磁准静态问题。
设两个互不相交的绝缘实心封装为 \(B_1,B_2\)，
\[
\Omega_\infty=\RR^3\setminus(\overline{B_1}\cup\overline{B_2}),
\qquad
\Omega_R=B_R(\bm c_0)\setminus(\overline{B_1}\cup\overline{B_2}).
\]
第一层是 \(\Omega_\infty\) 上满足无穷远物理衰减的场；
第二层是 \(\Omega_R\) 上指定了\emph{人工}外边界条件的连续问题；
第三层才是有限维试探函数及数值积分得到的离散问题。
\textbf{有限维功率闭合、有限外域相邻解接近、磁核矩阵半正定，
分别不能证明这三层问题已经一致。}

首版只假设静止、线性、非磁性绝缘封装、
\(\mu=\mu_0\)、均匀正海水电导率 \(\sigma_s\)、
闭合的外加 BFZI 两端口电流源及峰值相量 \(e^{\jj\omega t}\)。
不把本模型扩大解释为含磁芯非线性、导电屏蔽、
位移电流、运动感应和铜线内全耦合涡流的通用参考解。

\subsection{从 Maxwell 方程严格得到电流弱式}
在海水无外加体电流的区域，磁准静态方程为
\begin{align}
\Curl\bm H&=\bm J_s,&
\Curl\bm E&=-\jj\omega\bm B,&
\Div\bm B&=0,\\
\bm B&=\mu_0\bm H,&
\bm J_s&=\sigma_s\bm E,&
\Div\bm J_s&=0.
\end{align}
完整闭合源线电流的单位磁矢势记为
\(\bm A_{{\rm src},p}=\mu_0\bm C_p\)，
\(\bm C_p=\frac1{4\pi}\int_{\Gamma_p}\mathrm d\bm\ell'/|\bm x-\bm r'|\)。
取合适规范下的自由空间核
\((\mathcal G_0\bm J)(\bm x)=\int_{\Omega}\bm J(\bm y)/
(4\pi|\bm x-\bm y|)\,\mathrm dV_y\)，则
\[
\sigma_s^{-1}\bm J_s+\jj\omega\mu_0\mathcal G_0\bm J_s
+\grad\phi=-\jj\omega\mu_0\sum_p\bm C_p I_p.
\]
对有限外域，若使用人工绝缘截断条件，正确的测试空间为
\begin{equation}
\boxed{
\mathcal V_R=\{\bm v\in H(\operatorname{div};\Omega_R):
\Div\bm v=0,\quad
\bm v\cdot\bm n=0\ \text{于 }
\partial B_1\cup\partial B_2\cup\partial B_R\}.
}
\label{eq:g1-full-boundary-space}
\end{equation}
关键是\textbf{整个有限域边界}都要出现在零法向条件中；
否则
\[
\int_{\Omega_R}\overline{\bm v}\cdot\grad\phi\,\mathrm dV
=\int_{\partial\Omega_R}\phi\,\overline{\bm v}\cdot\bm n\,\mathrm dS
\]
在人工外球处未必为零，不能从弱式中凭空删除该边界项。
如采用非反射人工边界、允许非零外球电流通量，则必须
\textbf{保留并正确处理}对应边界算子/边界功项，
不能套用同一个无边界项弱式。
无限域则应使用具有相应远场衰减和磁能可积性的函数空间，
证明无穷远边界项消失后再作消元。

对所有 \(\bm v\in\mathcal V_R\) 得
\begin{equation}
a_R(\bm J_s,\bm v)
=\int_{\Omega_R}\sigma_s^{-1}\bm J_s\cdot\overline{\bm v}\,\mathrm dV
+\jj\omega\mu_0\int_{\Omega_R}
(\mathcal G_0\bm J_s)\cdot\overline{\bm v}\,\mathrm dV
=-\jj\omega\sum_p I_p\,(\bm A_{{\rm src},p},\bm v)_{L^2}.
\label{eq:g1-weak-corrected}
\end{equation}
这是\emph{人工绝缘有限域问题}的正确变分形式，
并不表示该问题等于无限海水的限制。

\subsection{哪些有限维结构结论确实正确}
选取实值线性独立的 \(\bm W_m\in\mathcal V_R\)，
以本稿式 \eqref{eq:Rs}--\eqref{eq:Kbfzi}
精确或一致求积定义 \(\bm R,\bm L,\bm K\)。
对实正 \(\sigma_s\)、自由空间对称 Coulomb 核及闭合端口，
理论上
\[
\bm R=\bm R^{\mathsf T}\succ0,\quad
\bm L=\bm L^{\mathsf T}\succeq0,\quad
(\bm R+\jj\omega\bm L)\bm z=-\jj\omega\bm K\bm I.
\]
令 \(\bm A=\bm R+\jj\omega\bm L\)，则
\[
\bm Z_{\rm sw}
=\omega^2\bm K^{\mathsf T}\bm A^{-1}\bm K,\qquad
\operatorname{Herm}(\bm Z_{\rm sw})
=\omega^2\bm K^{\mathsf T}\bm A^{-\mathrm H}
\bm R\bm A^{-1}\bm K\succeq0.
\]
故 \(\bm Z_{\rm sw}^{\mathsf T}=\bm Z_{\rm sw}\)，且
\[
\tfrac12\bm I^{\mathrm H}\operatorname{Herm}(\bm Z_{\rm sw})\bm I
=\tfrac12\bm z^{\mathrm H}\bm R\bm z
=\int_{\Omega_R}\frac{|\bm J_s|^2}{2\sigma_s}\,\mathrm dV.
\]
这些是\textbf{同一个离散系统的代数定理}；
它们证明约束、互易和耗散结构，不给出
\(\|\bm J_s-\bm J_\infty\|\) 或热点误差的任何界。
以离散自身的 \(\bm R,\bm K\) 重建功率，
即使 \(\bm L\) 或电流空间错误，仍可在机器精度闭合。

\subsection{当前 curl--高斯基的数学优势与缺口}
实现取 \(\bm W_m=\Curl\bm F_m\)，
\(\bm F_m\) 是高斯包络、绝缘盒体距离因子及外球截断因子
的乘积再乘常向量或完整端口的源磁矢势。
在适当的分片 \(H(\operatorname{curl})\) 正则性及相容迹下，
若 \(\bm n\times\bm F_m=0\)，可保证
\(\Div\bm W_m=0\) 与 \(\bm n\cdot\bm W_m=0\)；
这是本方案可保留的正确设计。

但当前外球 \(C^2\) 截断函数同时满足
\(f(R)=f'(R)=0\)，因此在外球处实际强制
\(\bm W_m|_{\partial B_R}=0\)，不仅是法向电流为零。
这在有限模态下是\emph{过于受限的合规子空间}，
并不等价于整个 \(\mathcal V_R\)；
它在 \(L^2\)/\(H(\operatorname{div})\) 完备性层面
未必构成不可逾越的阻碍，但需要\textbf{系统的空间加密证明}，
不能由 229 个高斯方向或有限几种长度尺度自行推出。
改变外球半径还会改变所有模式的边界包络，
故即使每个高斯中心、宽度及方向保持不变，
\textbf{跨半径函数空间也不是严格嵌套的}；
相邻解差异同时包含外域、外球边界和试探空间改变。

两个不相交、各自同胚于三维闭球的\emph{实心}绝缘盒体
在三维球域中的补集，一般没有一维同调环的拓扑障碍；
\textbf{不能由“有两个封装/多个边界分支”直接推断
必须加入非零的一阶拓扑谐和电流}。
物理上仍需表示跨封装及海水远场的\emph{普通大尺度涡流}，
这是空间逼近问题而不是自动出现的拓扑缺陷。
几何若改成穿孔实体、导电回路等，才需要重新计算拓扑条件。

\subsection{磁 Green 离散与连续 Coulomb 双积分不可混淆}
目标磁算子是
\[
L_{mn}
=\frac{\mu_0}{4\pi}
\int_{\Omega_R}\int_{\Omega_R}
\frac{\bm W_m(\bm x)\cdot\bm W_n(\bm y)}
{|\bm x-\bm y|}\,\mathrm dV_y\mathrm dV_x.
\]
目前 \texttt{uniform\_ball} 规则把每个求积权重视作等体积
\emph{球形电流云}，用球--球平均 Coulomb 核替代真实单元内的
双体积分。这能构造对称非负的近似磁能，
但球未必是实际正权积分单元的形状或支撑，
也未必保持真实海水边界和点值电流的无散度结构；
\textbf{正定性不蕴含对上述目标算子的数值一致性}。
必须独立检查采样节点分辨率、球形替代误差、
近对角奇异项与非局部尾项；当模式数和外域改变时不能
仅复用低精度 Green 矩阵并宣称参考级精度。
当前 229 模态 A/B \(\|\Delta L\|_F/\|L_A\|_F\approx13.2\%\)，
确实没有算子收敛证据。
这一 Frobenius 相对量本身也不能代替
\(\omega\|\bm R^{-1/2}\Delta\bm L\bm R^{-1/2}\|_2\)
或被激发方向的残差响应敏感度；
需同时检查这些\emph{无量纲动力学相关指标}，
不能仅由一个全矩阵百分比判断各端口场的误差。

\subsection{实际海水穿透深度与截断半径}
对于远离封装的均匀导电海水，在 MQS 假设下
\[
\grad^2\bm H=\jj\omega\mu_0\sigma_s\bm H,\qquad
\gamma=\sqrt{\jj\omega\mu_0\sigma_s}
=\frac{1+\jj}{\delta_s},\qquad
\delta_s=\sqrt{\frac{2}{\omega\mu_0\sigma_s}}.
\]
默认示例 \(f=200\,\mathrm{kHz}\)，
\(\sigma_s=5\,\mathrm{S/m}\)，
\(\mu=\mu_0\) 时 \(\delta_s\simeq0.503\,\mathrm m\)。
目前 \(0.52\)--\(1.30\,\mathrm m\) 的球半径仅约
\(1.03\)--\(2.58\) 个穿透深度（从球心计）；
因此外域远场电流及其磁反馈未必足够小，
\textbf{不能预设这些半径应形成收敛平台}。
球形/非均匀几何中的振幅还含角向及代数衰减，
实际截断误差不能只用 \(e^{-R/\delta_s}\) 代替测量。
较大的球域也会增加 Green 积分的求积误差：
外域相邻解非单调并不能单凭日志判断是物理尾部还是积分别名。

\subsection{从代数残差升级到 G1 可解释的误差界}
在固定\emph{连续}有限域问题中定义
\(\|\bm u\|_{\mathcal R}^2=
\int_{\Omega_R}\sigma_s^{-1}|\bm u|^2\,\mathrm dV\)。
因 \(\operatorname{Re}a_R(\bm u,\bm u)
=\|\bm u\|_{\mathcal R}^2\)，
对真实有限域解 \(\bm J_R\) 和任意合规近似
\(\bm J_h\)，其在\emph{完整测试空间}上的残差
\[
\mathfrak r_h(\bm v)
=\ell_R(\bm v)-a_R(\bm J_h,\bm v),\qquad
\|\mathfrak r_h\|_{\mathcal V_R'}
=\sup_{\bm v\in\mathcal V_R\setminus\{0\}}
\frac{|\mathfrak r_h(\bm v)|}{\|\bm v\|_{\mathcal R}}
\]
满足
\begin{equation}
\boxed{\|\bm J_R-\bm J_h\|_{\mathcal R}
\leq\|\mathfrak r_h\|_{\mathcal V_R'}.}
\label{eq:g1-residual-continuum}
\end{equation}
该结论是对\emph{相同连续算子与完整空间}的有条件估计，
不是当前离散 229 父空间上的残差上界。
只在另一个有限测试字典上取 supremum 时，
得到的是完整对偶残差范数的\emph{下界}，
除非再建立可靠的稳定性、饱和性或误差上界估计。
因此目前 \texttt{parent\_dJ\_energy\_upper}
只证明\emph{给定父空间内部}的逼近效果，
不能作无限维电磁参考误差认证。

对空间焦耳热有连续积分控制
\[
\|q(\bm J)-q(\bm J_h)\|_{L^1(\Omega_R)}
\leq\tfrac12\|\bm J-\bm J_h\|_{\mathcal R}
\big(\|\bm J\|_{\mathcal R}+\|\bm J_h\|_{\mathcal R}\big),
\]
但这并不控制盒壁极局部的 \(L^\infty\) 热点；
必须另外检查几何随体的壁面/边缘/角部点列及更强局部范数。

\subsection{第一批可实施的数学重构（小空间实验，而非直接认证）}
\paragraph{不修改 Coulomb 核的近场奇异分拆}
为避免将任意高斯求积权重伪装成\emph{球形体积单元}，
固定海水测试空间和有限域，取
\(\bm x\in\Omega_R\)，
\(0<a(\bm x)<\operatorname{dist}(\bm x,\partial\Omega_R)\)。
构造 \(C^2\) 径向窗 \(\chi_a(\rho)\)，
在 \(\rho\le a/2\) 时为 \(1\)，在 \(\rho\ge a\) 时为 \(0\)。
则对任何可积向量场 \(\bm W_n\) 有严格恒等式
\begin{align}
\int_{\Omega_R}\frac{\bm W_n(\bm y)}{|\bm x-\bm y|}\,\mathrm dV_y
={}&\int_{\Omega_R}
\frac{1-\chi_a(|\bm x-\bm y|)}
{|\bm x-\bm y|}\bm W_n(\bm y)\,\mathrm dV_y
\nonumber\\
&+\int_{S^2}\int_0^{a(\bm x)}
\rho\,\chi_a(\rho)\,
\bm W_n(\bm x+\rho\bm s)\,\mathrm d\rho\,\mathrm d\Omega_{\bm s}.
\label{eq:g1-coulomb-exact-split}
\end{align}
第二项在物理球坐标中解析消除了 \(1/\rho\) 奇异因子，
\emph{不引入核软化或节点云球模型}。
这是核恒等式，离散计算仍须分别对
\(\bm x\)、远场 \(\bm y\)、近场径向和角度作独立加密，
并保证相应求积的积分区域与模式评估精确一致。
实现 \texttt{sdfmpneo\_bfzi/green\_coulomb.py} 是有限模态小型
对照器，不能直接宣称其粗节点解精确或可快速运行 229 模态；
使用不同 \(\bm x/\bm y\) 求积时暂以对称化矩阵诊断，
\textbf{不能强迫其正定以掩盖积分误差}。

\paragraph{可逐级加密的紧支 curl 图谱}
用 \(C^2\) Wendland 核
\[
\psi(s)=(1-s)_+^4(4s+1),\qquad
\grad\psi(\sqrt{\bm r^{\mathsf T}G\bm r})
=-20(1-s)_+^3G\bm r
\]
替换局部高斯分布的固定有限个采样方向，
在封装外的面片以 \(2^\ell\times 2^\ell\)
网格状\emph{中心选点}生成紧支向量势，
再取 \(\Curl(\psi\,d_{\rm ins}\,\bm a)\)。
这不是有限元体网格，而是\emph{解析紧支基函数的层级索引}，
可在层级加密时保持所有已存在的电流方向与
精确解析 curl。额外建立按物理封装长度而非人工球半径
选择中心的远场图谱。实现
\texttt{build\_sea\_basis(structured\_level=1,2)} 已构造此嵌套字典；
目前尚无空间稠密性定理或跨几何误差界。

\paragraph{真正无限域的解析边界基准}
把任意 BFZI 双盒问题直接附加一个经验 Robin 常数是错误的。
先在最简单、可解析的绝缘球腔内中心磁偶极模型中求解。
对偶极的轴对称环向模 \(A_\phi=f(r)\sin\theta\)，
令 \(\gamma^2=\jj\omega\mu_0\sigma_s\)、
\(\operatorname{Re}\gamma>0\)，则海水中
\[
f''+\frac2r f'-\frac2{r^2}f-\gamma^2 f=0,\qquad
f_\infty(r)=Ck_1(\gamma r),
\]
\(k_1\) 为衰减的修正球 Bessel 函数。
在任何包围球腔的半径 \(R\)，此\emph{唯一球谐扇区}
的精确无限域透明条件为
\begin{equation}
\boxed{\frac{f_\infty'(R)}{f_\infty(R)}
=\gamma\frac{k_1'(\gamma R)}{k_1(\gamma R)}.}
\label{eq:g1-exact-dtn-l1}
\end{equation}
内侧真空向量势由 \(M/r^2+Dr\) 组成；
在腔壁连续匹配 \(f\) 和 \(f'\)，
可以得到不含人工外球的精确海水 \(J_\phi=-\jj\omega\sigma_s
f_\infty\sin\theta\) 与 \(q=|J_\phi|^2/(2\sigma_s)\)。
这已在 \texttt{sdfmpneo\_bfzi/exterior\_dipole.py} 中实现，
同时可比较把外球的 \(A_\phi\) 强制为零导致的截断误差。
\textbf{式 \eqref{eq:g1-exact-dtn-l1} 不适用于任意 BFZI
线圈、多极混合和两个长方体封装。}
推广须处理完整矢量球谐的极向/环向模及耦合边界映射，
或改用严谨的边界积分施加非局部透明条件。

\subsection{真实封装几何裁剪的固定半径 Coulomb 奇异补偿}
第一版为保证近场球全部落在海水，使用
\(a(\bm x)=\min(a_0,0.6\,\mathrm{dist}(\bm x,\partial\Omega_R))\)。
在封装面附近，这会使 \(a(\bm x)\to0\)；
理论恒等式并未失效，但需要由粗全局 \(\bm y\) 求积
解析越来越尖锐的近对角核，易出现积分次序别名与相消。
现在允许固定物理补偿半径 \(a_0>0\)，
但\emph{严格按真实海水域裁剪}球形补偿区：
\[
\int_{\Omega_R\cap B_{a_0}(\bm x)}
\frac{\chi_{a_0}(|\bm y-\bm x|)\bm W_n(\bm y)}
{|\bm y-\bm x|}\,\mathrm dV_y
=\int_{S^2}\sum_{[s_k,t_k]\subset I_{\bm x,\bm n}}
\int_{s_k}^{t_k}\rho\,\chi_{a_0}(\rho)
\bm W_n(\bm x+\rho\bm n)\,\mathrm d\rho\,\mathrm d\Omega,
\]
其中 \(I_{\bm x,\bm n}\) 是沿固定射线
\(\bm x+\rho\bm n\) 的海水区间与 \([0,a_0]\) 的交，
用真实旋转盒的射线--板交点和外球交点构造，
不在封装内部评估任何电流模态。
全局 \(\bm y\) 部分仍使用
\((1-\chi_{a_0})/|\bm x-\bm y|\)，
因此\emph{连续}双积分依旧严格等于原始 Coulomb 算子。

此修正仅移除了\emph{近壁补偿球被迫收缩}的结构性限制，
并未使数值求积自动收敛。因为球被盒体切割，
沿方向 \(\bm n\) 的积分具有遮挡/切触轮廓，
角向采样也必须独立加密；必须同时检查
\(a_0\)、局部角向阶数、\(\bm x/\bm y\) 体积规则，
以及对驱动相关的 \(\omega R^{-1/2}LR^{-1/2}\) 的影响。
2026-11-04 的旧规则实测
\(dL_{\bm x}=9.8049\%\)、\(dL_{\bm y}=3.4771\%\)、
\(dL_{\bm x\bm y}=2.9804\%\)，
非加性指标 \(10.4896\%\)，
不能仅据双轴误差抵消宣称新旧任何一方已收敛。

\paragraph{针对近壁方向非光滑性的局部求积修正（新代码待实测）}
2026-11-04 的真实几何裁剪试验已经证明，
近壁补偿球不能只依赖均匀角向规则：
近区角向 \(4/8\to6/12\) 的 \(dL\) 从旧缩球的
\(4.51\times 10^{-7}\) 增大到 \(3.60734\times10^{-3}\)，
改变补偿半径 \(a:0.055\to0.080\,\mathrm m\)
使 \(dL=2.11519\%\)。
独立加密双重积分的外层 \(\bm x\) 给出
\(dL_{\bm x}=11.2368\%\)，内层 \(\bm y\) 仅
\(1.37432\%\)；两轴同时变化 \(1.14953\%\)
来自明显的非加性，并非物理精度通过。

改进方法不改变原 Coulomb 核、实际海水区域、
基函数或局部补偿球：
对最邻近的一处盒壁/人工球面选择局部球面极轴
\(\bm n\)，并将角向积分划分为两段 Gauss 面板。
距无限平面盒壁正距离 \(d<a\) 时，第一处半径
\(a\) 射线触壁对应 \(\mu_*=\cos\theta_*=d/a\)，
取面向固体的法向为极轴。
对于中心 \(\bm c_0\)、半径 \(R\) 的人工外球，
\[
\mu_*=\frac{R^2-|\bm x-\bm c_0|^2-a^2}
{2a|\bm x-\bm c_0|},
\]
极轴沿球心指向 \(\bm x\)。
两个面板 \([-1,\mu_*]\) 和 \([\mu_*,1]\)
各自使用独立的 Gauss--Legendre 极角积分，
方位角使用周期求积；
截断盒角/盒边以及多盒遮挡仍由射线交点计算，
此角向坐标选择\emph{不改变}真实湿区积分。
与此同时，在径向窗 \(\chi_a\) 的分片连接点
\(\rho=a/2\) 显式切开每一个湿区区间，
减少单个低阶 Gauss 规则穿越 \(C^2\) 分片连接
时的高阶导数非光滑性。

平面绝缘壁常向量场有可复核的独立基准：
令采样点距壁 \(d\)，局部球半径 \(a>d\)，则
\[
\int_{\Omega_{\rm half-space}\cap B_a(\bm x)}
\frac{\chi_a(|\bm y-\bm x|)}
{|\bm y-\bm x|}\,\mathrm dV_y
=\int_0^d 4\pi\rho\chi_a(\rho)\,\mathrm d\rho
+\int_d^a 2\pi(\rho+d)\chi_a(\rho)\,\mathrm d\rho .
\]
此积分在固定角向规则下可以直接比较旧的单极角
与新的面向实体的分段极角，且不是人工制造的
原核软化或拟合标签。
它\emph{仅}验证局部半空间的角向裁剪；
G1 仍要求真实旋转双盒场景的
\(\bm x/\bm y\) 求积、分割半径、近壁局部热及
驱动方向稳定性，并检查不同 \(\ell_{\max}\)
的远区物理匹配。

\subsection{均匀导电海水的完整球谐外区映射：当前实现边界}
选择包围所有绝缘封装、完整端口源的球面
\(S_R\)，球外只有均匀正标量电导率海水。
由 \(\Div\bm J=0\)、
\(\Curl\bm E=-\jj\omega\mu_0\bm H\)、
\(\Curl\bm H=\bm J=\sigma_s\bm E\) 得
\[
(\Delta-\gamma^2)\bm J=0,\qquad
\gamma^2=\jj\omega\mu_0\sigma_s,\quad\Re\gamma>0.
\]
故在固定 Cartesian 分量下，对\emph{散度相容}的边界迹
\(\bm J(R,\theta,\varphi)\)，外区衰减解为
\begin{equation}
J_i(r,\theta,\varphi)
=\sum_{\ell=0}^{\infty}\sum_{m=-\ell}^{\ell}
j_{i,\ell m}(R)
\frac{k_\ell(\gamma r)}{k_\ell(\gamma R)}
Y_{\ell m}(\theta,\varphi),\qquad r\ge R.
\end{equation}
球面透明算子是角向非局部的
\begin{equation}
\partial_r J_i(R,\theta,\varphi)
=\sum_{\ell,m}
\underbrace{\gamma\frac{k'_\ell(\gamma R)}
{k_\ell(\gamma R)}}_{\lambda_\ell(R)}
j_{i,\ell m}(R)Y_{\ell m}(\theta,\varphi).
\label{eq:g1-multipolar-dtn}
\end{equation}
这解释了为什么\textbf{只有一个} Robin 系数在一般
双盒 BFZI 问题中不正确，且其 \(\ell=1\)
子块与前述解析偶极基准相符。
利用球谐正交性，已知相容球面电流系数时，外海水
平均 Joule 功率为
\[
P_{\mathrm{tail}}
=\frac1{2\sigma_s}\sum_{i,\ell,m}
|j_{i,\ell m}(R)|^2
\int_R^\infty r^2\left|
\frac{k_\ell(\gamma r)}{k_\ell(\gamma R)}
\right|^2\,\mathrm dr .
\]

实现 \texttt{ExteriorSphericalCurrent} 提供
有限 \(\ell_{\max}\) 截断下的投影、外区续接、
复对数导数与尾区损耗，
并加入 \(\ell=1,2\) 环向无散度模态数值基准。
\textbf{关键未完成：}
从有限域原始 Coulomb Galerkin 解得到的任意球面迹
未必满足无限域散度相容条件，且现有外球包络将
边界电流强制压零。必须将式
\eqref{eq:g1-multipolar-dtn}
作为非局部边界条件\emph{反向耦合}到双盒内区求解，
并控制角向截断、匹配球半径与真物理边界误差。
只对某个已截断的数值场做球谐外推，
不构成完整 BFZI 无穷域误差证书。

\subsection{纠正后的 G1 通过条件与工作次序}
\begin{enumerate}
\item \textbf{正确边值问题：}
证明每个频率与材料参数下源回路闭合，
说明无限域衰减与人工外域边界算子的关系，
区分非导电封装的零法向和人为外球条件；
检查 MQS 适用范围（例如 \(\omega\epsilon/\sigma\ll1\)）。
\item \textbf{真实磁核：}
在固定物理基及同一个外域上，独立加密 Green 积分、
控制对角近奇异项、球云形状替代和远场尾项，
给出动力学加权磁算子及双端口受迫场的一致误差，
而非仅验证 PSD。
\item \textbf{有系统的海水空间逼近：}
建立随尺度加密的紧支 curl/Piola 近盒、间隙、
完整远场角向/径向字典，保证合规与适当稠密性；
对独立大父空间评估 J、双端口 J、q 和残差，
不是只在固定 229 父解里做压缩。
\item \textbf{外域独立收敛：}
固定共同物理湿区，同时用真正物理一致的
外域尾场/透明边界或多个充分大的截断半径核对；
必须把外域截断、边界包络变化、R/K 求积和 Green L
积分误差作为\emph{不同}来源逐项控制。
\item \textbf{独立物理基准：}
至少在代表性几何、错位姿态和封装尺寸下与
不同离散原则的高精度解（解析可用极限、
独立边界积分或其他可信参考）交叉核对，
报告端口、海水体积 J、封装近壁 q 和功率误差。
允许参考方法使用局部离散；
主代理仍可坚持无固定世界网格设计。
\end{enumerate}
在以上各项证据成立以前，
\textbf{G1=未通过，229 模态不能用作正式跨几何神经教师}。
下一轮优先分离同一物理坐标中的 R/K 与 Green L
敏感性，接着修磁核求积和边界/空间逼近；
不再把简单扩大外域或增加少数经验高斯模式
当作充分的 G1 收敛方案。
